% Methodology（草稿：数据/特征/模型/SCS/UI/风控/评估协议）
\label{sec:method}

\subsection{Data}
\label{sec:data}
Figure~\ref{fig:framework} gives an overview of the full pipeline; the
components are detailed below. The evaluation uses 11,811 league matches from the English Premier League (E0),
Spanish La Liga (SP1), German Bundesliga (D1), Italian Serie A (I1),
and French Ligue 1 (F1), covering the 2019/20--2025/26 seasons
(2025/26 is incomplete, ending 2026-02-12; this is reported honestly).
Each match record contains results, referee, and opening and closing
odds from eight bookmakers (Bet365, Bet\&Win, Interwetten, Pinnacle,
William Hill, Ladbrokes, Max, and average) in football-data.co.uk
format. Per-match expected goals (xG, xGA,
npxG, npxGA), pressing intensity (PPDA), and deep progression are
additionally crawled from Understat for all five leagues and all seven seasons
(24,918 team-match rows; 99.3\%/99.2\% of home/away sides match
league-match dates after name normalization).

\paragraph{Leakage control.}
The evaluation is strictly chronological:
train = 2019/20--2023/24 (8,955 matches), validation = 2024/25
(1,752), test = 2025/26 (1,104). All features are constructed from
information available before kickoff; rolling statistics use a
shift of one to exclude the current match; season standings are
computed before kickoff; scaler statistics and missing-value medians
are fit on the training split only and applied to validation and
test. A manual verification script recomputes a sample of features
from raw records to confirm the absence of leakage.

\subsection{Features}
All features are pre-match:
\begin{itemize}
\item \textbf{Team form:} rolling (5-match) means of goals for/
against, points, win rate, shots, shots on target, corners, fouls;
season cumulative points, goal difference, games, points per game,
and league rank before kickoff (home and away separately).
\item \textbf{xG time series:} rolling (5-match) means of xG, xGA,
npxG, npxGA, PPDA coefficient, deep progressions, and expected
points, plus the home--away xG difference.
\item \textbf{Market signals:} de-vigged implied probabilities from
closing odds (Bet365 and average), opening-to-closing odds movement
for home/draw/away, cross-bookmaker dispersion of closing odds
($\sigma_{market}$), and opening average odds.
\item \textbf{Referee context:} the referee's historical per-match
averages of cards, red cards, and fouls (expanding window over the
referee's prior matches; available for the Premier League only).
\end{itemize}
This yields 60 features in total.

\subsection{Models}
\begin{itemize}
\item \textbf{Market baseline:} de-vigged closing odds
(probabilities obtained by inverting odds and renormalizing). This
is the central baseline under the same evaluation protocol.
\item \textbf{Poisson} and \textbf{Dixon--Coles} models fit per
league by maximum likelihood on the training split.
\item \textbf{Elo} ratings with home advantage, mapped to
probabilities by multinomial logistic regression fit on training.
\item \textbf{Random Forest} and \textbf{XGBoost} trained on the 60
pre-match features with early stopping on the validation split.
\item \textbf{Domain-enhanced LLM:} DeepSeek-Chat, prompted with a
pre-match feature card (team form, rank, market signals, referee
context), a system prompt stating that closing odds are a strong
baseline and forbidding invented information, and a chain-of-thought
request that outputs a calibrated probability triple in JSON. The
LLM has no retrieval access. Three samples are drawn per match at
temperature 0.3; the mean is the prediction and the pairwise
total-variation agreement is the consistency score.
\end{itemize}

\subsection{Scenario Complexity Score}
\label{sec:scs}
The scenario complexity score measures how hard a match is to
predict from \emph{market information alone}. Unlike earlier
formulations in which a larger strength gap was treated as
\emph{more} complex (a direction error), the corrected SCS is
defined as
\[
\mathrm{SCS} = (1 - \mathrm{gap}_{norm}) + \mathrm{vol}_{norm}
+ \mathrm{move}_{norm},
\]
where $\mathrm{gap}_{norm}$ is the normalized absolute league-rank
difference between home and away teams (larger gap $\Rightarrow$
easier to predict), $\mathrm{vol}_{norm}$ is the robust-normalized
cross-bookmaker dispersion of closing odds, and
$\mathrm{move}_{norm}$ is the robust-normalized opening-to-closing
odds movement (a proxy for late information arrival). All components
lie in $[0,1]$ and are fit on the training split. A larger SCS
indicates a more uncertain match.

\subsection{Uncertainty Index and Risk Control}
\label{sec:ui}
The uncertainty index is conceptually distinct from the SCS: SCS is a
\emph{purely market-side} measure of scenario complexity (rank gap,
odds dispersion, odds drift), whereas the UI \emph{fuses model output
with market signals}---predictive confidence and multi-sample
consistency from the model, plus the same normalized market-volatility
and odds-movement terms. The shared market terms are intentional: the
UI inherits market uncertainty and adds the model's own confidence
signal. The index combines these components as
\[
\mathrm{UI} = w_1(1 - p_{\max}) + w_2(1 - C)
+ w_3\,\mathrm{vol}_{norm} + w_4\,\mathrm{move}_{norm},
\]
\noindent with default weights $(0.4, 0.3, 0.15, 0.15)$; $p_{\max}$ is the
model's maximum predicted probability and $C$ the multi-sample
consistency (1 for deterministic models). The consistency term is
retained even though Section~\ref{sec:llm_analysis} shows it carries
no information in this setting: it is the mechanism proposed by
earlier LLM-forecasting work, and including it lets this study test
and falsify that mechanism rather than silently dropping it; the
weight sensitivity analysis confirms the UI's stratification is
insensitive to this term. Risk tiers are
$\mathrm{UI} \le t_{low}$ (full stake), $t_{low} < \mathrm{UI} \le
t_{high}$ (linearly scaled stake), and $\mathrm{UI} > t_{high}$
(no-bet). Thresholds are selected on the validation split
($t_{low}=0.30$, $t_{high}=0.45$ given the observed UI distribution)
and reported with a sensitivity analysis. A stop-loss rule reduces
the stake to 20\% after five consecutive losses or a 10\% drawdown
from peak bankroll.

\subsection{Betting Protocol and Metrics}
\label{sec:protocol}
All financial simulations are fully specified: equal stakes of one
unit (or fractional Kelly capped at 10\% of bankroll where stated),
the model's maximum-probability outcome, Bet365 closing odds for
settlement, sequential bankroll updating, and no-bet decisions
excluded from executed-bet statistics but reported separately with
coverage. Predictive metrics are accuracy, macro-F1, log loss,
Brier score, and expected calibration error (ECE). Financial metrics
are ROI, Sharpe ratio, maximum drawdown, and win rate. All main
quantities are reported with bootstrap 95\% confidence intervals.
